Let $X=\Leb^\infty(0,\pi)$.  The paper compares the construction with four sufficient conditions
for maximal $\Leb^2$-regularity at the endpoint and deduces their failure from
Theorem~\ref{thm:main}: if one of them held, the cited theorem would give maximal regularity.
The cited theorems are not formalised.  The formalisation instead proves the failure of every
hypothesis directly from the construction, which does not depend on the literature.

\begin{lemma}[The operator norm of a form increment]
  \label{lem:opnorm}
  \lean{LionsC12.formOpDist, LionsC12.enorm_form_sub_le, LionsC12.formOpDist_le_eLpNorm,
    LionsC12.exists_form_sub_ge_of_lt_eLpNorm, LionsC12.formOpDist_eq_eLpNorm,
    LionsC12.formOpDist_a, LionsC12.aX, LionsC12.norm_aX_sub}
  \leanok
  \uses{def:lions-solution, def:coefficient}
  For a bounded measurable coefficient on an interval $(a,b)$ with continuous increments,
  $\norm{\cA(t)-\cA(s)}_{\cL(V,V')}=\norm{\fra(t)-\fra(s)}_{\Leb^\infty(a,b)}$.
\end{lemma}

\textbf{Lean form.}  The paper uses only the inequality $\leq$.  The operator norm is
\texttt{formOpDist}, the supremum of
$\abs{\fra(t)[v,w]-\fra(s)[v,w]}/(\norm{v}_V\norm{w}_V)$ over $v,w\neq0$, with the graph norm of
$V$; no operator $V\to V'$ is constructed.  The path $t\mapsto\fra(t)\in X$ is \texttt{aX}, with
values in Mathlib's $\Leb^\infty$.

\begin{proof}
  \leanok
  The upper bound is H\"older's inequality.  For the lower bound, test with
  $v=\psi((x-x_0)/h)$, a dilated smooth bump at a point where the increment is close to its
  supremum, for which $\norm{\partial_xv}^2/\norm{v}_V^2\to1$ as $h\to0$.
\end{proof}

\begin{corollary}[Failure of the known sufficient hypotheses at the endpoint]
  \label{cor:endpoint}
  \lean{LionsC12.endpoint_hypotheses_fail, LionsC12.eVariationOn_aX_eq_top,
    LionsC12.not_exists_dier_bound, LionsC12.not_exists_dini_form,
    LionsC12.not_exists_dini_coeff, LionsC12.dini_form_lintegral_eq_top,
    LionsC12.dini_coeff_lintegral_eq_top, LionsC12.gagliardoCoeff, LionsC12.gagliardoForm,
    LionsC12.gagliardoForm_eq, LionsC12.gagliardoCoeff_eq_top,
    LionsC12.exists_piece_gagliardo_eq_top, LionsC12.aeSquare,
    LionsC12.not_exists_auscherEgert_bound, LionsC12.not_exists_auscherEgert_bound',
    LionsC12.iSup_lintegral_formOpDist_sq_div_le, LionsC12.exists_partition_local_smallness,
    LionsC12.form_self_re_equiv}
  \leanok
  \uses{lem:opnorm, prop:coefficient, prop:holder-sharp}
  For the coefficient of Theorem~\ref{thm:main}:
  \begin{enumerate}
    \item the form does not have bounded variation in the sense of Dier, and
      $\fra\notin BV([0,1];X)$;
    \item there is no non-decreasing $\omega:[0,1]\to[0,\infty)$ with
      $\abs{\langle(\cA(t)-\cA(s))v,w\rangle}\leq\omega(\abs{t-s})\norm{v}_V\norm{w}_V$ for
      $s,t\in[0,1]$ and $\int_0^1\omega(r)r^{-3/2}\dd r<\infty$ (Haak--Ouhabaz); in particular no
      such modulus controls $\norm{\fra(t)-\fra(s)}_X$;
    \item neither $\cA:[0,1]\to\cL(V,V')$ nor $\fra:[0,1]\to X$ is piecewise $\Sob^{1/2}$ on any
      finite partition (Achache--Ouhabaz);
    \item there is no $M<\infty$ with
      $\sup_{I}\abs{I}^{-1}\iint_{I\times I}\abs{\fra(t,x)-\fra(s,x)}^2/\abs{t-s}^2\dd s\dd t\leq M$
      for a.e.\ $x\in(0,\pi)$ (Auscher--Egert, condition (8)).
  \end{enumerate}
  Moreover the two side conditions of Achache--Ouhabaz hold: the local smallness
  $\sup_{t\in I}\int_I\norm{\cA(t)-\cA(s)}^2_{\cL(V,V')}\abs{t-s}^{-1}\dd s\leq C\eps^2\abs{I}$, and the
  uniform equivalence $\re\fra(t)[v,v]\simeq\norm{v}_V^2$.
\end{corollary}

\textbf{Lean form.}  The statement is \texttt{endpoint\_hypotheses\_fail}, for every
$0<\eps\leq1$.  (i) is \texttt{eVariationOn} of the path \texttt{aX} on $[0,1]$, equal to
$\infty$, and the absence of a function $g$, monotone on $[0,1]$, with Dier's estimate.  (ii)
reads ``$\int_0^1\omega(r)r^{-3/2}\dd r<\infty$'' as finiteness of the lower Lebesgue integral of
$\omega(r)r^{-3/2}$ on $(0,1)$; the stronger versions
\texttt{dini\_form\_lintegral\_eq\_top}, \texttt{dini\_coeff\_lintegral\_eq\_top} say that the
integral is $\infty$ for every admissible monotone $\omega$.  (iii) says that every partition
$0=t_0<\dots<t_N=1$ has a piece $(t_k,t_{k+1})$ on which the Gagliardo seminorm
$\iint\norm{\fra(t)-\fra(s)}_X^2\abs{t-s}^{-2}$ (and the same for $\cA$, Lemma~\ref{lem:opnorm}) is
infinite; only the seminorm of the $\Sob^{1/2}$ norm is used, and the piece is the one containing
the accumulation time $5/16$.  (iv) has the supremum over open intervals $(p,q)\subseteq[0,1]$
inside the almost-everywhere quantifier, as in the paper; \texttt{not\_exists\_auscherEgert\_bound'}
excludes even the weaker hypothesis in which the exceptional null set may depend on the interval.  The local smallness constant is
$C=4C_\beta^2$.  The form equivalence is stated with $\norm{\partial_xv}_H$ in place of
$\norm{v}_V$ (the two are equivalent by Poincar\'e's inequality, which is not formalised for
$V$; the upper bound in the graph norm is \texttt{norm\_form\_a\_le}); the operator square root of
the Kato property is not formalised.  The theorems of Dier, Haak--Ouhabaz, Achache--Ouhabaz and
Auscher--Egert are not formalised.

\begin{proof}
  \leanok
  This is the argument Lean states, not the paper's, which deduces (i)--(iv) from the cited
  theorems and Theorem~\ref{thm:main}.  On the middle third $J_j$ of a block, the increment of
  $\fra(\cdot,0)$ over a quarter period of the first mode is at least $\frac{13}{15}\eps K_{j,1}^{-1}$;
  $3K_{j,1}$ such steps fit in $J_j$, so each block contributes at least $\eps$ to the variation,
  which proves (i), and Dier's estimate would bound it.  For (ii), on the pieces
  $(h_{j+1},h_j]$ of the scales $h_j=\pi/(2K_{j,1}^2)$, monotonicity gives
  $\omega(r)r^{-3/2}\geq\omega(h_{j+1})h_j^{-3/2}$, each piece contributes at least $c/(j+2)$, and the
  harmonic series diverges.  For (iii) and (iv), on rectangles $T\times S\subseteq J_j\times J_j$
  where the phase of the first mode lies in $[\pi/4,3\pi/4]$ and in $[5\pi/4,7\pi/4]$, the
  squared difference quotient of $\sin K^2(\cdot-\tau_j)$ integrates to at least
  $K^2\abs{J_j}/(36\pi)$, which, after the factor $\eps^2K^{-2}$, is $\geq c\eps^2\ell_j$ per block
  on a set of length $\ell_j$, hence unbounded after division by $\abs{I}$ and infinite in sum; (iv)
  integrates over $x$ by Tonelli.  Local smallness is
  Proposition~\ref{prop:coefficient} with Lemma~\ref{lem:opnorm}.
\end{proof}
